\documentclass[exercice,a4]{thedocument}%
\startdatakeys
    \source{laclassecooperativedestef.com}
    \BO{2026}
    \cycle{3}
    \competencesDuSocle{RE4,CO2}
    \objectifsApprentissage{C3C1ab,C3C1aj,C3C2ab,C3C1ae,C3C1af,C3C2an,C3C1ag}
    \motsCles{milieu, perpendicularite, parallelisme, triangle}
\enddatakeys
\begin{document}%
\noindent\textbf{Programme de construction 15}\par\smallskip\noindent
\begin{itemize}
    \item Tracer le segment $[\mathspoint{AB}]$ tel que
    $\mathspoint{AB}=10\text{ cm}$.
    \item Placer $\mathspoint{C}$, milieu de $[\mathspoint{AB}]$.
    \item Placer $\mathspoint{D}$, milieu de $[\mathspoint{AC}]$.
    \item Placer $\mathspoint{E}$, milieu de $[\mathspoint{CB}]$.
    \item Tracer la droite $(\mathspoint{d})$ passant par
    $\mathspoint{C}$, perpendiculaire à $[\mathspoint{AB}]$.
    \item Placer $\mathspoint{F}$ sachant que $\mathspoint{CF}=5\text{ cm}$.
    \item Tracer les triangles $\mathspoint{DFE}$ et $\mathspoint{AFB}$.
    \item Tracer la droite $(\mathspoint{e})$ passant par
    $\mathspoint{F}$, parallèle à $[\mathspoint{AB}]$.
    \item Tracer la droite $(\mathspoint{f})$ passant par
    $\mathspoint{A}$, perpendiculaire à $(\mathspoint{e})$ ; elle coupe
    $(\mathspoint{e})$ en $\mathspoint{G}$.
    \item Tracer la droite $(\mathspoint{g})$ passant par
    $\mathspoint{B}$, perpendiculaire à $(\mathspoint{e})$ ; elle coupe
    $(\mathspoint{e})$ en $\mathspoint{H}$.
\end{itemize}
\end{document}%
